% !TeX program = xelatex
\documentclass[11pt,a4paper]{ctexart}

\usepackage{amsmath,amssymb,bm}
\usepackage{booktabs,longtable,array,multirow}
\usepackage{geometry}
\geometry{margin=2.4cm}
\usepackage{xcolor}
\usepackage{listings}
\usepackage{enumitem}
\usepackage[colorlinks=true,linkcolor=blue!60!black,urlcolor=blue!60!black]{hyperref}

\lstdefinestyle{code}{basicstyle=\ttfamily\footnotesize,frame=single,framerule=0.3pt,
  rulecolor=\color{black!40},backgroundcolor=\color{black!4},breaklines=true,
  columns=fullflexible,keepspaces=true,showstringspaces=false}
\lstset{style=code}

\newcommand{\file}[1]{\texttt{#1}}
\newcommand{\func}[1]{\texttt{#1}}
\newcommand{\opt}[1]{\texttt{#1}}
\newcommand{\evid}[1]{{\footnotesize\color{black!70}（证据：#1）}}

\title{\textbf{MIPSolvers 原生 IPM 体系增强路线图\\——2023--2026 前沿算法调研与并行化失效诊断}}
\author{基于 \texttt{docs/native\_ipm\_audit\_report.pdf} 审计结论的续篇}
\date{2026 年 8 月 2 日}

\begin{document}
\maketitle

\begin{abstract}
\noindent 本报告解决一个具体问题：\textbf{为什么当前并行计算不能为该 IPM 体系加速，以及如何用 2023--2026 年的最新算法成果修复它}。诊断结论：OpenMP 仅包裹了装配级 $O(\mathrm{nnz})$ 循环，而占每迭代 70--90\% 成本的\textbf{稀疏分解全部串行}（EigenSparseLU/KLU/UMFPACK 顺序执行、MUMPS 走 mpiseq 顺序兼容层、CHOLMOD 未配置并行）——按 Amdahl 定律，现状并行收益上限仅约 1.1--1.4$\times$。增强路线的核心判断：本代码库与当前 GPU-IPM 前沿（HyKKT/LiftedKKT 凝聚 KKT + cuDSS）的接口距离\textbf{比表面看起来近得多}——凝聚形式 $W=H+J_h^{\top}(M/S)J_h$、SPD 认证、CHOLMOD 封装、稠密 Schur 求解、``analyze-once/factorize-many''缓存契约均已存在；真正的缺口是（i）分解后端串行且不定、（ii）无真实惯性、（iii）审计发现的两个串行热点（线搜索全 Jacobian 重评、Merit 路径每迭代 $\le 4$ 次重复分解）封死了任何并行化的收益上限。
\end{abstract}

\tableofcontents

\section{并行化失效诊断：为什么现在并行不加速}

\subsection{当前并行触点全清单}
全库 OpenMP 在 IPM 内核中的使用点（编译开关 \opt{MIPSOLVERS\_USE\_OPENMP} 默认 ON，\file{CMakeLists.txt:439}）：
\begin{center}\footnotesize
\begin{tabular}{@{}lll@{}}
\toprule
位置 & 并行内容 & 量级 \\
\midrule
\file{cones.cpp:324} & NT 缩放的逐锥块计算 & $O(\sum q + \sum p^3)$，小 \\
\file{conic\_ipm\_solver.cpp:521,553,724} & Schur（$G^{\top}H^{-1}G$）装配 & $O(\mathrm{nnz})$ 级 \\
\file{ipm\_lp\_solver.cpp:1761} 等 & $\Theta$ 构造、残差合并等逐分量循环 & $O(n)$ \\
\bottomrule
\end{tabular}
\end{center}
线程数控制：锥求解器 \opt{num\_threads} 默认 0（库默认），\func{OmpThreadsGuard} 仅在 $>0$ 时设定\evid{\file{conic\_ipm\_solver.hpp:31}，\file{conic\_ipm\_solver.cpp:65--81}}。

\subsection{串行瓶颈：分解后端全部顺序执行}
每迭代主导成本是 KKT/正规方程分解，而所有默认后端均为\textbf{串行}：
\begin{center}\footnotesize
\begin{tabular}{@{}llll@{}}
\toprule
后端 & 类型 & 并行性 & 位置 \\
\midrule
EigenSparseLU & 非对称 LU & 串行 & \file{linear\_solver.cpp:26--58} \\
KLU/UMFPACK/SuperLU & 非对称 LU & 串行 & \file{linear\_solver.cpp:60--192} \\
MKL PardisoLU & 非对称 LU & 仅 MKL 内部线程（且为非对称路径） & \file{linear\_solver.cpp:194--237} \\
MUMPS & 对称不定 LDL$^{\top}$ & \textbf{mpiseq 顺序层，MPI\_Init 为空操作} & \file{linear\_solver.cpp:248--257} \\
CHOLMOD（锥 KKT） & SPD LL$^{\top}$/LDL$^{\top}$ & 未配置 GPU/多线程，仅超节点 GEMM 借 BLAS 线程 & \file{cholmod\_ldlt.cpp:37--58} \\
\bottomrule
\end{tabular}
\end{center}
注意 MUMPS 即便编译了并行，当前也通过 mpiseq 以单进程方式驱动\evid{\file{linear\_solver.cpp:248--257} 注释}。于是\textbf{分解这一主导环节零并行}。

\subsection{Amdahl 定量估计}
IPM 每迭代成本近似 $T = T_{fact} + T_{asm} + T_{eval} + T_{ls}$。对中大型稀疏问题，直接法分解 $T_{fact}$ 通常占 70--90\%（Lp 正规方程 $O(\mathrm{nnz}(N)^{1.5})$ 级，NLP 增广 KKT 更高）。当前并行只作用于 $T_{asm}$ 的一部分（且装配本身已有高效的 scatter-map 缓存，占比进一步压低），故
\begin{equation}
S_{max} \le \frac{1}{(1-p)} \approx 1.1\text{--}1.4\times \quad (p \le 0.3 \text{ 的可并行份额，且其中仅一部分被实际并行}),
\end{equation}
与``并行无加速''的观察一致。此外审计发现两个\textbf{串行热点}会进一步封顶任何并行收益，必须先修：
\begin{itemize}[leftmargin=2em]
  \item \textbf{H1}：NLP-IPM 线搜索每次回退重算全部 Jacobian（\file{ipm\_solver.cpp:849--884} $\to$ 736--753）——Ipopt 只评 $f,g,h$；
  \item \textbf{H2}：Merit 路径 reduced-KKT 每迭代 $\le 4$ 次对同一 $W$ 全量重新分解（\file{ipm\_solver.cpp:2317,2357,2408} $\to$ \file{kkt\_system.cpp:283--354}）。
\end{itemize}

\section{前沿调研（2023--2026）：与本体系相关的最新 IPM 算法}

\subsection{GPU 凝聚 KKT 内点法：HyKKT / LiftedKKT / MadNLP}
Pacaud、Shin、Montoison、Schanen、Anitescu 的系列工作是当前 GPU-NLP-IPM 的主线\footnote{\href{https://arxiv.org/abs/2405.14236}{Pacaud et al., \emph{Condensed-space methods for nonlinear programming on GPUs}, arXiv:2405.14236 (2024)}。}：
\begin{itemize}[leftmargin=2em]
  \item \textbf{核心观察}：KKT 直接分解的数值主元（pivoting）本质串行，是 GPU 化的根本障碍；将不定 KKT \textbf{凝聚为 SPD 系统}后可用无主元 Cholesky，天然适合 GPU。
  \item \textbf{HyKKT}：将 $(1,1)$ 块（$H_\gamma = H + $ 正则，SPD）用 GPU Cholesky 分解，等式 Schur 补 $S_\gamma = J_g H_\gamma^{-1} J_g^{\top}$ 用 \textbf{共轭梯度（CG）}迭代求解——完全无主元；分析表明凝聚系统的病态是``结构化的''，精度损失可控。
  \item \textbf{LiftedKKT}：不等式提升为等式+松弛后凝聚为单一 SPD 系统，直接 Cholesky。
  \item 实现在 MadNLP.jl（filter 线搜索 IPM，与本体系同框架）+ ExaModels.jl（GPU 导数评估）；经典蒸馏塔算例\textbf{每迭代提速 25 倍}\footnote{\href{https://arxiv.org/abs/2403.15913}{Shin et al., \emph{GPU-accelerated dynamic nonlinear optimization with ExaModels and MadNLP}, arXiv:2403.15913 (2024)}。}；已用于 67.4 万维电网问题。
\end{itemize}

\subsection{cuDSS：GPU 稀疏直接求解器}
NVIDIA cuDSS（2024）提供 GPU 驻留的稀疏直接分解；其契约是\textbf{分析阶段（METIS 排序，CPU）一次、数值分解在 GPU 反复重用}\footnote{\href{https://discourse.julialang.org/t/sparse-lu-factorization-on-gpu/122012}{A. Montoison, CUDSS.jl 讨论 (2024)}；\href{https://docs.nvidia.com/cuda/cudss/}{NVIDIA cuDSS 文档}。}——与本代码库 \func{SparseKKTCache} 的``结构指纹缓存 + analyze once / factorize many''契约\evid{\file{kkt\_system.cpp:74--148,206--242}}\textbf{完全同构}，移植阻抗极低。

\subsection{CPU 侧：DAG 任务并行的对称不定分解}
不引入 GPU 的现实路线。HSL MA86（OpenMP、DAG 任务调度、对称不定）\footnote{\href{https://www.hsl.rl.ac.uk/specs/hsl_ma86.pdf}{HSL MA86 规范：\emph{Sparse symmetric indefinite system using OpenMP}。}}、HSL MA87/MA97（DAG 型，MA86 的不定对应物）、SPRAL/SSIDS（CPU+GPU 混合，且\textbf{能提供惯性}——HSL 邮件列表明确指出 Ipopt 场景下只有 SPRAL 类能提供所需惯性信息\footnote{\href{https://list.coin-or.org/pipermail/ipopt/2016-August/004295.html}{COIN-OR Ipopt 列表：Sparse Linear Solver on GPU for IPOPT（J. Hogg）}。}）。Ipopt 官方支持 \opt{ma86}/\opt{ma97}/\opt{spral} 后端\footnote{\href{https://coin-or.github.io/Ipopt/OPTIONS.html}{Ipopt 官方选项文档（linear\_solver 一节）}。}。GALAHAD SLS 提供统一接口封装 MA57/MA86/MA97/SSIDS/MUMPS/PARDISO/PaStiX\footnote{\href{https://ralna.github.io/galahad_docs/html/C/sls.html}{GALAHAD SLS 文档}。}——\textbf{这正是本库 \func{SparseLinearSolver} 接口应该长成样子}。

\subsection{一阶 GPU LP 方法：PDLP $\to$ cuPDLP $\to$ cuPDLP+}
对超大 LP（特别是 MIPLIB 松弛），无分解的一阶方法已成第三极：cuPDLP（重启 PDHG，仅稀疏 matvec，GPU 全并行）\footnote{\href{https://arxiv.org/pdf/2506.02174}{Lu \& Yang, \emph{An Overview of GPU-based First-Order Methods for LP}, arXiv:2506.02174 (2025)。}}；cuPDLP+ 引入 Halpern 重启+反射、新重启判据、PID 原始权重更新，在 MIPLIB 松弛上再快 2--4$\times$\footnote{\href{https://arxiv.org/abs/2507.14051}{Lu, Peng \& Yang, \emph{cuPDLP+}, arXiv:2507.14051 (2025)。}}。局限同样明确：精度有限（$10^{-4}$ 级实用）、无基解、暖启动能力弱——\textbf{在 B\&C 内部不能替代 IPM/单纯形，但可作为超大独立 LP 的前置粗解器}。本库已有 \func{NativePDLPAdapter}（\file{src/engine/solver/native/lp/pdlp\_solver.cpp}），升级路径现成。

\subsection{批量与结构化并行}
\begin{itemize}[leftmargin=2em]
  \item \textbf{Many-problems-one-GPU}：将大量中小 NLP 批量驻留单 GPU 并行求解\footnote{\href{https://arxiv.org/html/2606.26341v1}{\emph{Scaling Nonlinear Optimization: Many Problems One GPU}, arXiv:2606.26341 (2026)。}}——直接对应 B\&C 强分支/泵浦中\textbf{成批节点 LP}的场景（本库 \func{solve\_cached\_bound\_change\_batch} 当前为串行循环，\file{ipm\_lp\_solver.cpp:3392}）。
  \item \textbf{块三对角 Schur 分解}（动态/NMPC 结构）：自动检测块结构 + GPU batched LAPACK，性能对比 MA57/CHOLMOD/cuDSS 有优势\footnote{\href{https://proceedings.aiche.org/conferences/aiche-annual-meeting/2025/proceeding/paper/gpu-accelerated-solution-block}{AIChE 2025: \emph{GPU-Accelerated Solution of Block-Tridiagonal Systems for Large-Scale Optimization}。}}——对应 SCUC/OPF 的时段耦合结构（本库 SCUC 元数据完备）。
  \item \textbf{PIPS-IPM/PIPS-NLP}：块角结构问题的 Schur 补分解 + MPI 并行，是结构化大规模 IPM 的经典范式（SCUC 多场景/多时段可直接套用）。
  \item \textbf{QOCO-GPU}：GPU 加速的二次目标 SOCP 内点法\footnote{\href{https://arxiv.org/html/2603.29197v1}{\emph{QOCO-GPU: A Quadratic Objective Conic Optimizer with GPU Acceleration}, arXiv:2603.29197 (2026)。}}——证明锥 IPM（本库最强内核）GPU 化路线同样成立；其 KKT $K=G^{\top}H^{-1}G+\delta I$ 为 SPD，与 cuDSS/CHOLMOD-GPU 直接兼容。
\end{itemize}

\section{前沿 $\to$ 本代码库映射}

\begin{center}\footnotesize
\begin{longtable}{@{}p{3.4cm}p{5.2cm}p{5.2cm}@{}}
\toprule
前沿 & 本库已有基础（锚点） & 缺口 \\
\midrule
\endhead
HyKKT 凝聚+CG & 凝聚 $W=H+J_h^{\top}(M/S)J_h$ 装配（\file{ipm\_solver.cpp:1260--1267}）；SPD 认证（\file{kkt\_system.cpp:384--407}）；小 Schur 稠密求解（\file{kkt\_system.cpp:283--354}） & 无 SPD 后端复用（每次新建分解器）；无 CG/Krylov 求解器；Schur 仅用于 $m_{eq}\le 8$ \\
cuDSS GPU 分解 & analyze-once 缓存契约（\file{kkt\_system.cpp:206--242}）；锥 KKT 为 SPD（CHOLMOD 路径） & 无 GPU 后端；值数组布局需适配 cuDSS 接口 \\
MA86/97、SPRAL 并行不定分解 & 后端抽象 \func{SparseLinearSolver}（\file{linear\_solver.hpp:12--19}） & 默认链全为非对称串行 LU；MUMPS 被 mpiseq 串行化 \\
真实惯性驱动 $\delta_W$ & MUMPS INFOG(12) 已封装 \func{negative\_eigenvalues()}（\file{linear\_solver.cpp:397--401}） & MUMPS 非默认；现有``惯性''是 $Z^{\top}WZ$ 稠密特征解（nullity$\le 128$）或全空间凸化 \\
cuPDLP+ 一阶 LP & \func{NativePDLPAdapter} 已存在（\file{pdlp\_solver.cpp}） & 无 Halpern 重启/PID 权重/GPU matvec \\
批量节点 LP & 批量接口存在但串行（\file{ipm\_lp\_solver.cpp:3392--3497}） & 无批量分解/batched BLAS \\
SCUC 块结构 & UCGenHint 元数据完备（\file{problem\_types.hpp:281--406}） & IPM 未利用块角/块三对角结构 \\
\bottomrule
\end{longtable}
\end{center}

\section{分层增强路线图}

\subsection{T1：CPU 快速收益（1--2 周，低风险）}
\begin{enumerate}[leftmargin=2em]
  \item \textbf{T1.1 修复串行热点（前置，否则一切并行收益被封顶）}：H1——线搜索试点只评 $f,g,h$（Jacobian 每迭代一次，试算 $\theta,\varphi$ 只需函数值）；H2——Merit 路径跨 RHS 复用同一 $W$ 的分解（predictor/corrector/Gondzio 共用一次分解）。预期 1.2--2$\times$（串行）。
  \item \textbf{T1.2 接入并行对称不定后端}：将 \func{make\_default\_sparse\_solver} 的 NLP-KKT 默认改为（按可得性）MKL PARDISO 对称不定路径（$LDL^{\top}$，多线程）、多线程 MUMPS（脱离 mpiseq，配多线程 BLAS；OpenMP 版 MUMPS）、或 SSIDS/SPRAL。预期分解环节 2--6$\times$（8--16 核）。
  \item \textbf{T1.3 真实惯性驱动 $\delta_W$}：以 MUMPS/SSIDS 的负主元计数（\func{negative\_eigenvalues()} 已封装）替换 $Z^{\top}WZ$ 稠密特征解路径，δ\_W 循环回到 Wächter--Biegler 标准形式；同时消除 nullity$>128$ 时的全空间凸化副作用。
  \item \textbf{T1.4 稠密路径换多线程 BLAS/LAPACK}：LCQP 与 LP-IPM 的稠密/带状 Cholesky（自研 \func{banded\_chol\_*}）换 LAPACK \func{dpbsv/dpotrf}（Accelerate/MKL 均已多线程）；稠密正规方程 $AA^{\top}$ 用 \func{dsyrk}。
\end{enumerate}

\subsection{T2：结构级（1--2 月，中风险）}
\begin{enumerate}[leftmargin=2em]
  \item \textbf{T2.1 HyKKT 化（无主元化，为 GPU 铺路）}：在现有凝聚路径上，将 $W+\delta_W I$（已有 SPD 认证逻辑）用 CholmodLDLT \textbf{一次分解、多 RHS 复用}；等式 Schur 补 $J_g(W{+}\delta I)^{-1}J_g^{\top}+\delta_C I$ 从``$m_{eq}\le 8$ 稠密''推广为任意 $m_{eq}$ 的 CG（matvec = 两次三角回代 + $J_g$ 稀疏乘，全部 GPU 友好）。收益：KKT 全流程无主元 $\Rightarrow$ 可直接落 cuDSS；并消除非对称 LU 的 2$\times$ 开销。\textbf{这是性价比最高的一步}。
  \item \textbf{T2.2 SCUC 块结构 Schur 分解}：利用 UCGenHint 的时段结构，将 KKT 重排为块角/块三对角形，对角块独立分解（OpenMP 任务并行）+ 耦合 Schur 串行小系统（PIPS-IPM 范式；动态结构可参考 AIChE-2025 块三对角 GPU 方案）。
  \item \textbf{T2.3 B\&C 批量节点 LP}：\func{solve\_cached\_bound\_change\_batch} 改为多线程分派（不同 bound-change 独立求解，线程安全依赖每线程 CachedState）；强分支候选 LP 同理批量化。
\end{enumerate}

\subsection{T3：GPU 化（1 季度以上，需 NVIDIA GPU + CUDA 工具链）}
\begin{enumerate}[leftmargin=2em]
  \item \textbf{T3.1 cuDSS 后端}：实现 \func{SparseLinearSolver} 的 cuDSS 版本（CPU/METIS 分析一次 $\to$ GPU 反复分解/回代），首先服务锥 IPM 的 SPD KKT（改动最小、收益最直接）。
  \item \textbf{T3.2 凝聚 SPD + 全 GPU KKT}：T2.1 的凝聚路径整体落 GPU（cuDSS Cholesky + GPU CG）；导数矩阵经现有 scatter-map 在 GPU 上填充值数组。参照 MadNLP 的全驻留设计，避免每迭代 H2D 拷贝。
  \item \textbf{T3.3 NativePDLP 升级}：引入 cuPDLP+ 的 Halpern 重启、反射、PID 原始权重；matvec 移 GPU。定位：超大独立 LP 的粗解器，不替代 B\&C 内 IPM。
  \item \textbf{T3.4 GPU 导数评估}：现回调式 \func{NLPModel} 无法 GPU 化；长期需表达式图（参照 ExaModels 的 SIMD/SIMT 单遍评估）。此项成本最高，建议仅在有明确大规模 NLP 业务需求时立项。
\end{enumerate}

\section{预期收益与风险（诚实估计）}
\begin{center}\footnotesize
\begin{tabular}{@{}llll@{}}
\toprule
路线 & 预期加速 & 主要风险 & 备注 \\
\midrule
T1.1 热点修复 & 1.2--2$\times$（串行） & 低 & 必须先做 \\
T1.2 并行不定后端 & 分解 2--6$\times$ & 中（HSL 许可/MKL 依赖） & 端到端受 Amdahl 限制 \\
T1.3 真实惯性 & 稳健性为主 & 低 & 附带去掉稠密特征解成本 \\
T2.1 HyKKT 化 & 1.5--3$\times$ + GPU 通行证 & 中（Schur-CG 收敛依赖 $\delta_C$） & 需 IR 兜底（已有 2 步 IR） \\
T2.2 结构分解 & 视块数 2--10$\times$ & 高（重排/映射复杂） & 仅结构化问题 \\
T2.3 批量节点 LP & 近线程数线性 & 中（线程安全） & B\&C 侧收益 \\
T3 GPU & 大问题 5--25$\times$/迭代 & 高（平台、精度、惯性缺失） & 分解 $>0.1$ s/次才划算（HSL 经验） \\
\bottomrule
\end{tabular}
\end{center}
通用告诫：(i) 凝聚路径失去分解惯性，$\delta_W$ 策略需重新设计（HyKKT 论文用固定正则+IR，本库的约化空间证书恰可补位）；(ii) GPU 仅在单次分解足够大时回本；(iii) 一阶方法精度上限决定其不能进 B\&C 内环。

\section{验证计划}
\begin{enumerate}[leftmargin=2em]
  \item \textbf{正确性门槛}：\file{tests/test\_ipm\_solver.cpp}、\file{tests/test\_conic\_ipm.cpp} 全绿；新增 KKT 残差 $\|Kz-b\|_\infty$ 与 KKT 误差回归测试（每后端）。
  \item \textbf{性能基准}：netlib/MIPLIB LP 松弛（LP-IPM）、CUTEst 中选 10 个中小 + 2 个大 NLP（NLP-IPM）、SDPLIB 若干（锥）；记录每迭代分解/装配/线搜索分项耗时（代码已有 verbose 计时基建）。
  \item \textbf{并行扩展性}：1/2/4/8/16 线程扫描，报告分解环节与端到端两组加速比（现状只有后者无意义）。
\end{enumerate}

\end{document}
